%%%%%%%%%%%%%%%%%%%%%%%%%%%%%%%%%%%%%%%%%
% Long Lined Cover Letter
% LaTeX Template
% Version 2.0 (September 17, 2021)
%
% This template originates from:
% https://www.LaTeXTemplates.com
%
% Authors:
% Vel (vel@LaTeXTemplates.com)
%
% License:
% CC BY-NC-SA 4.0 (https://creativecommons.org/licenses/by-nc-sa/4.0/)
%
%%%%%%%%%%%%%%%%%%%%%%%%%%%%%%%%%%%%%%%%%

%----------------------------------------------------------------------------------------
%	PACKAGES AND OTHER DOCUMENT CONFIGURATIONS
%----------------------------------------------------------------------------------------

\documentclass[10pt]{article}

\usepackage{charter} % Use the Charter font

\usepackage[
	a4paper, % Paper size
	top=.5in, % Top margin
	bottom=.7in, % Bottom margin
	left=1in, % Left margin
	right=1in, % Right margin
	%showframe % Uncomment to show frames around the margins for debugging purposes
]{geometry}

\setlength{\parindent}{0pt} % Paragraph indentation
\setlength{\parskip}{1em} % Vertical space between paragraphs

\usepackage{graphicx} % Required for including images
\usepackage{natbib}
\usepackage[colorlinks=true, linkcolor=magenta, citecolor=magenta, urlcolor=magenta]{hyperref}
\usepackage{fancyhdr} % Required for customizing headers and footers

\fancypagestyle{firstpage}{%
	\fancyhf{} % Clear default headers/footers
	\renewcommand{\headrulewidth}{0pt} % No header rule
	\renewcommand{\footrulewidth}{1pt} % Footer rule thickness
}

\fancypagestyle{subsequentpages}{%
	\fancyhf{} % Clear default headers/footers
	\renewcommand{\headrulewidth}{1pt} % Header rule thickness
	\renewcommand{\footrulewidth}{1pt} % Footer rule thickness
}

\AtBeginDocument{\thispagestyle{firstpage}} % Use the first page headers/footers style on the first page
\pagestyle{subsequentpages} % Use the subsequent pages headers/footers style on subsequent pages

%----------------------------------------------------------------------------------------

\begin{document}

%----------------------------------------------------------------------------------------
%	FIRST PAGE HEADER
%----------------------------------------------------------------------------------------

%\includegraphics[width=0.35\textwidth]{logo.pdf} % Logo

\vspace{-1em} % Pull the rule closer to the logo

\rule{\linewidth}{1pt} % Horizontal rule

\bigskip\bigskip % Vertical whitespace

%----------------------------------------------------------------------------------------
%	YOUR NAME AND CONTACT INFORMATION
%----------------------------------------------------------------------------------------

\hfill
\begin{tabular}{l @{}}
	\today \bigskip\\ % Date
\href{https://mamintoosi.github.io/}{Mahmood Amintoosi}\\
Division of Computer Science,\\
Department of Applied Mathematics \&\\
\href{https://gta-lab.github.io/}{Graph Theory and its Applications Lab} \\
Faculty of Mathematical Sciences,\\ Ferdowsi University of Mashhad,
IRAN
%\url{https://mamintoosi.github.io/}%{https://mamintoosi.github.io/}
\end{tabular}

%\bigskip % Vertical whitespace

%----------------------------------------------------------------------------------------
%	ADDRESSEE AND GREETING
%----------------------------------------------------------------------------------------

%\begin{tabular}{@{} l}
%	Mrs.\ Jane Smith \\
%	Recruitment Officer \\
%	The Corporation \\
%	123 Pleasant Lane \\
%	City, State 12345
%\end{tabular}

\bigskip % Vertical whitespace

Dear Editor,

\bigskip % Vertical whitespace

%----------------------------------------------------------------------------------------
%	LETTER CONTENT
%----------------------------------------------------------------------------------------
I am pleased to submit our original research article, ``Graph Structure Learning for Traffic Prediction'', for consideration for publication in the \textit{International Journal of Data Science and Analytics}.
Traditional graph-based methods for traffic forecasting rely on predefined adjacency matrices based on physical road proximity, which often fail to capture true functional dependencies between roads. In this work, we introduce a knowledge discovery framework that learns the underlying graph structure directly from traffic data using continuous optimization, shifting from heuristic graph construction to data-driven knowledge extraction. Key contributions are as follows:
\begin{itemize}
    \item A novel approach that infers meaningful, persistent traffic dependencies directly from data, overcoming the limitations of predefined spatial relationships.
    \item Two methodological variants: GSL (directed acyclic graph capturing causal temporal influences) and cGSL (symmetrized cyclic version for spatial models).
    \item Comprehensive validation on two real-world datasets (SZ-Taxi, Los-loop) demonstrating significant improvements—up to 24.7\% RMSE reduction for GCN-cGSL and 21.8\% for TGCN-GSL.
\end{itemize}

This work aligns well with the journal's scope, advancing both traffic prediction and graph-based knowledge discovery. It shares common ground with recent contributions such as \cite{Chen2026TASTGCN} on trend-aware spatio-temporal graph convolutional networks.
While a preliminary version was presented at a national conference \citep{Amintoosi1403aimc55-GSL}, this submission represents a substantial extension:
\begin{itemize}
    \item \textbf{Scope:} Extends both GCN and T-GCN architectures (conference focused only on T-GCN).
    \item \textbf{Implementation:} Complete modern implementation in PyTorch 2.0 with Google Colab notebooks for reproducibility  available at \url{https://github.com/mamintoosi/TGCN-GSL-PyTorch} (conference used deprecated TensorFlow 1.x).
    \item \textbf{Experiments:} Comprehensive evaluation on two datasets across multiple horizons (conference presented only preliminary results).
    \item \textbf{Novelty:} Introduces the cyclic cGSL variant and provides detailed ablation studies absent from the conference version.
\end{itemize}

%The conference paper is properly cited and discussed within the manuscript. All code and data are publicly available at \url{https://github.com/mamintoosi/TGCN-GSL-PyTorch}.

%We believe this work will be of great interest to your readers and look forward to your editorial decision.

%Thank you for considering my paper for publication in ``Data \& Knowledge Engineering". I look forward to the opportunity to contribute to the journal.


\bigskip % Vertical whitespace

Sincerely,

%\vspace{50pt} % Vertical whitespace

Mahmood Amintoosi

\bibliographystyle{elsarticle-harv} 
\bibliography{MyReferences}

\end{document}
